\documentclass[border=5pt]{standalone}
\usepackage[utf8]{inputenc}
\usepackage[spanish]{babel}
\usepackage{amsmath,amssymb}
\usepackage{tikz}
\usetikzlibrary{positioning,arrows.meta,fit,calc,shapes.geometric}
\begin{document}
\definecolor{cEnt}{RGB}{31,78,121}
\definecolor{cSem}{RGB}{31,110,60}
\definecolor{cEnc}{RGB}{101,60,140}
\definecolor{cFus}{RGB}{176,32,42}
\definecolor{cSal}{RGB}{191,110,20}
\begin{tikzpicture}[
  font=\scriptsize,
  mod/.style={rounded corners=2.5pt, draw, semithick, align=center, inner sep=4pt, anchor=north},
  tarj/.style={rounded corners=2.5pt, draw=black!25, fill=black!2, semithick,
               align=left, inner sep=5pt, text width=3.6cm, anchor=north},
  chico/.style={font=\tiny, text=black!60},
  fl/.style={-{Latex[length=1.9mm]}, semithick, black!60}
]
%% ---------- rama semantica (arriba) ----------
\node[mod, draw=cSem, fill=cSem!5, text width=3.1cm] (M2) at (5.6,16.2)
 {\textcolor{cSem}{\textbf{2. ¿A qué se parece cada punto?}}\\[1pt]
  {\tiny\itshape\color{cSem!85}rama semántica de HyperNCD}\\[3pt]
  \tikz[baseline=0, every node/.style={anchor=center}]{
    \fill[cSem!25, draw=cSem] (0,0.35) -- (0.5,0.16) -- (0.5,-0.16) -- (0,-0.35) -- cycle;
    \fill[cSem!45, draw=cSem] (0.68,0.16) -- (1.18,0.35) -- (1.18,-0.35) -- (0.68,-0.16) -- cycle;
    \draw[-{Latex[length=1.4mm]}, black!45] (0.5,0) -- (0.68,0);}\\[4pt]
  Una red entrenada con clases ya etiquetadas dice a qué pieza conocida se
  parece: muro, columna, arco.\\[3pt]
  \tikz[baseline=0, every node/.style={anchor=center}]{
    \foreach \i in {0,...,4} \fill[cSem!30, draw=cSem!70] (\i*0.14,0) rectangle (\i*0.14+0.12,0.15);
    \node[font=\tiny, anchor=west] at (0.72,0.075) {$\cdots$};}\\[1pt]
  {\tiny\color{black!70}$f_{\mathrm{sem}}$: la lista de 96 números que esa red
  le pone a cada punto para resumir \textbf{a qué se parece}}};

%% ---------- rama auto-supervisada (abajo) ----------
\node[mod, draw=cEnc, fill=cEnc!5, text width=3.1cm] (M3) at ($(M2.south)+(0,-0.55)$)
 {\textcolor{cEnc}{\textbf{3. ¿Qué forma hay alrededor?}}\\[1pt]
  {\tiny\itshape\color{cEnc!85}codificador de Sonata}\\[3pt]
  \tikz[baseline=0, every node/.style={anchor=center}]{
    \foreach \i/\w in {0/1.0, 1/0.84, 2/0.68, 3/0.52, 4/0.36}
      \fill[cEnc!35, draw=cEnc] (0,\i*0.15) rectangle (\w,\i*0.15+0.11);
    \foreach \a in {0,60,120} \draw[cyan!65!black, semithick] (1.28,0.33) ++({\a}:0.15) -- ++({\a+180}:0.30);}\\[4pt]
  Otra red, preentrenada sola y sin etiquetas, describe la forma del entorno.
  Se usa congelada.\\[3pt]
  \tikz[baseline=0, every node/.style={anchor=center}]{
    \foreach \i in {0,...,4} \fill[cEnc!30, draw=cEnc!70] (\i*0.14,0) rectangle (\i*0.14+0.12,0.15);
    \node[font=\tiny, anchor=west] at (0.72,0.075) {$\cdots$};}\\[1pt]
  {\tiny\color{black!70}$f_{\mathrm{ssl}}$: la lista de 1088 números que esa red
  le pone a cada punto para resumir \textbf{cómo es la forma}}};

\coordinate (eje) at ($(M2.east)!0.5!(M3.east)$);

%% ---------- entrada ----------
\node[mod, draw=cEnt, fill=cEnt!5, text width=2.3cm, anchor=east] (M1) at ($(eje)-(4.85,0)$)
 {\textcolor{cEnt}{\textbf{1. Lo que entra}}\\[1pt]
  {\tiny\itshape\color{cEnt!85}datos del proyecto}\\[3pt]
  \tikz[baseline=0, every node/.style={anchor=center}]{\foreach \x/\y in {0.214/-0.013, 0.225/0.095, 0.221/0.215, 0.204/0.330, 0.203/0.438, 0.205/0.535, 0.217/0.672, 0.206/0.760, 0.365/0.016, 0.363/0.106, 0.377/0.204, 0.373/0.322, 0.347/0.426, 0.353/0.561, 0.349/0.663, 0.365/0.765, 1.442/-0.016, 1.424/0.099, 1.446/0.217, 1.433/0.333, 1.438/0.433, 1.451/0.557, 1.431/0.663, 1.441/0.784, 1.588/-0.008, 1.597/0.096, 1.577/0.229, 1.567/0.330, 1.563/0.446, 1.590/0.553, 1.594/0.653, 1.587/0.773, 1.583/0.778, 1.579/0.929, 1.527/1.046, 1.450/1.165, 1.386/1.279, 1.289/1.338, 1.156/1.414, 1.015/1.446, 0.888/1.446, 0.751/1.457, 0.626/1.399, 0.518/1.359, 0.404/1.259, 0.336/1.172, 0.283/1.053, 0.225/0.910, 0.215/0.794, 1.456/0.767, 1.418/0.876, 1.389/0.986, 1.352/1.071, 1.264/1.159, 1.195/1.231, 1.123/1.286, 1.006/1.314, 0.906/1.304, 0.809/1.320, 0.707/1.290, 0.596/1.225, 0.504/1.167, 0.435/1.064, 0.391/0.974, 0.365/0.869, 0.342/0.767, 0.086/-0.065, 0.198/-0.047, 0.334/-0.073, 0.436/-0.065, 0.555/-0.074, 0.688/-0.042, 0.789/-0.061, 0.890/-0.074, 1.014/-0.068, 1.147/-0.072, 1.233/-0.044, 1.366/-0.073, 1.482/-0.077, 1.596/-0.043, 1.723/-0.053, 1.021/1.055, 0.988/1.153, 0.924/1.198, 0.829/1.163, 0.789/1.122, 0.791/1.027, 0.846/0.956, 0.913/0.933, 0.994/0.959} \fill[cEnt!75] (\x,\y) circle (0.019);}\\[4pt]
  Nube de puntos del escaneo láser: la posición y el color de cada punto.\\[3pt]
  {\tiny\color{black!60}$N$ puntos · vóxel de 5\,cm}};

%% ---------- fusion ----------
\node[mod, draw=cFus, very thick, fill=cFus!4, text width=2.6cm, anchor=west] (M4) at ($(eje)+(0.9,0)$)
 {\textcolor{cFus}{\textbf{4. Se juntan y se agrupa}}\\[1pt]
  {\tiny\itshape\color{cFus!85}razonamiento de HyperNCD}\\[3pt]
  \tikz[baseline=0, every node/.style={anchor=center}]{
    \fill[cSem!45, draw=cSem] (0,0) rectangle (0.42,0.22);
    \fill[cEnc!45, draw=cEnc] (0.42,0) rectangle (2.1,0.22);}\\[4pt]
  Cada punto lleva las dos descripciones a la vez; los puntos parecidos se
  reúnen en grupos y los grupos se comparan entre sí.};

%% ---------- salida ----------
\node[mod, draw=cSal, fill=cSal!6, text width=2.2cm, anchor=west] (M5) at ($(M4.east)+(1.5,0)$)
 {\textcolor{cSal}{\textbf{5. Lo que sale}}\\[1pt]
  {\tiny\itshape\color{cSal!85}resultado del sistema}\\[3pt]
  \tikz[baseline=0, every node/.style={anchor=center}]{\foreach \x/\y in {0.214/-0.013, 0.225/0.095, 0.221/0.215, 0.204/0.330, 0.203/0.438, 0.205/0.535, 0.217/0.672, 0.206/0.760, 0.365/0.016, 0.363/0.106, 0.377/0.204, 0.373/0.322, 0.347/0.426, 0.353/0.561, 0.349/0.663, 0.365/0.765, 1.442/-0.016, 1.424/0.099, 1.446/0.217, 1.433/0.333, 1.438/0.433, 1.451/0.557, 1.431/0.663, 1.441/0.784, 1.588/-0.008, 1.597/0.096, 1.577/0.229, 1.567/0.330, 1.563/0.446, 1.590/0.553, 1.594/0.653, 1.587/0.773, 1.583/0.778, 1.579/0.929, 1.527/1.046, 1.450/1.165, 1.386/1.279, 1.289/1.338, 1.156/1.414, 1.015/1.446, 0.888/1.446, 0.751/1.457, 0.626/1.399, 0.518/1.359, 0.404/1.259, 0.336/1.172, 0.283/1.053, 0.225/0.910, 0.215/0.794, 1.456/0.767, 1.418/0.876, 1.389/0.986, 1.352/1.071, 1.264/1.159, 1.195/1.231, 1.123/1.286, 1.006/1.314, 0.906/1.304, 0.809/1.320, 0.707/1.290, 0.596/1.225, 0.504/1.167, 0.435/1.064, 0.391/0.974, 0.365/0.869, 0.342/0.767, 0.086/-0.065, 0.198/-0.047, 0.334/-0.073, 0.436/-0.065, 0.555/-0.074, 0.688/-0.042, 0.789/-0.061, 0.890/-0.074, 1.014/-0.068, 1.147/-0.072, 1.233/-0.044, 1.366/-0.073, 1.482/-0.077, 1.596/-0.043, 1.723/-0.053, 1.021/1.055, 0.988/1.153, 0.924/1.198, 0.829/1.163, 0.789/1.122, 0.791/1.027, 0.846/0.956, 0.913/0.933, 0.994/0.959} \fill[cSal!80] (\x,\y) circle (0.019);}\\[4pt]
  Cada punto con su clase, y lo nunca visto agrupado aparte.\\[3pt]
  {\tiny\color{black!60}mIoU \emph{known} / \emph{novel}}};

%% ---------- flechas ----------
\draw[fl] (M1.east) to[out=22, in=180, looseness=0.9] (M2.west);
\draw[fl] (M1.east) to[out=-22, in=180, looseness=0.9] (M3.west);
\draw[fl] (M2.east) -- (M4.west);
\draw[fl] (M3.east) -- (M4.west);
\draw[fl] (M4.east) -- (M5.west);
\node[chico, align=center, fill=white, inner sep=1.5pt]
  at ($(M1.east)!0.46!(M2.west)$) {la misma\\nube};
\node[chico, fill=white, inner sep=1pt] at ($(M2.east)!0.45!(M4.west)+(0,0.16)$) {$f_{\mathrm{sem}}$};
\node[chico, fill=white, inner sep=1pt] at ($(M3.east)!0.45!(M4.west)+(0,-0.16)$) {$f_{\mathrm{ssl}}$};
\node[chico, anchor=south, fill=white, inner sep=1.5pt]
  at ($(M4.east)!0.5!(M5.west)+(0,0.1)$) {etiquetas};

%% ================= DETALLE DEL MODULO 4 =================
\coordinate (pc) at ($(M3.south)+(1.7,-0.85)$);
\node[anchor=north west, font=\bfseries, text=cFus] (tdet) at ($(pc)+(-7.4,0)$)
  {Dentro del Módulo 4: qué se calcula y cómo se lee};

\node[tarj] (T1) at ($(pc)+(-5.45,-1.15)$)
 {\textbf{\textcircled{\scriptsize 1} Se juntan las dos descripciones}\\[4pt]
  \centering
  \tikz[baseline=0, every node/.style={anchor=center}]{
    \fill[cSem!45, draw=cSem] (0,0) rectangle (0.42,0.24);
    \fill[cEnc!45, draw=cEnc] (0.42,0) rectangle (2.6,0.24);
    \node[font=\tiny] at (0.21,0.12) {96};
    \node[font=\tiny] at (1.51,0.12) {1088};
    \draw[{Latex[length=1.2mm]}-{Latex[length=1.2mm]}, black!45] (0,-0.13) -- (2.6,-0.13);
    \node[font=\tiny, text=black!60] at (1.3,-0.35) {1184 números};
    \fill[cSem!45, draw=cSem] (0,-1.05) rectangle (0.42,-0.81);
    \fill[red!30, draw=cFus] (0.42,-1.05) rectangle (0.55,-0.81);
    \node[font=\tiny, text=black!60, anchor=west] at (0.64,-0.93) {línea base: 96+6};}\\[6pt]
  \centering$F_i=\left[\,f_{\mathrm{sem}};\;f_{\mathrm{ssl}}\,\right]$\\[6pt]
  \raggedright
  El corchete significa pegar una lista de números detrás de la otra: la ficha
  del punto $i$ pasa a tener $96+1088=1184$ valores. Antes esa segunda parte
  eran 6 números calculados a mano, y \textbf{cambiarlos es el aporte}.};

\node[tarj] (T2) at ($(pc)+(0,-1.15)$)
 {\textbf{\textcircled{\scriptsize 2} Cada punto se reparte entre los grupos}\\[4pt]
  \centering
  \tikz[baseline=0, every node/.style={anchor=center}]{
    \foreach \x/\y/\c in {0.222/0.339/cSem, 0.457/0.245/cSem, 0.212/0.293/cSem, 0.401/0.416/cSem, 0.234/0.177/cSem, 0.429/0.353/cSem, 0.494/0.392/cSem, 0.330/0.247/cSem, 1.343/0.558/cEnc, 1.069/0.463/cEnc, 1.179/0.624/cEnc, 1.092/0.570/cEnc, 1.182/0.660/cEnc, 1.267/0.602/cEnc, 0.986/0.644/cEnc, 1.005/0.562/cEnc, 1.651/0.120/cSal, 1.712/0.144/cSal, 1.999/0.117/cSal, 1.884/-0.012/cSal, 1.766/0.197/cSal, 1.785/0.179/cSal, 1.811/0.011/cSal, 1.862/0.031/cSal} \fill[\c!55] (\x,\y) circle (0.035);
    \foreach \x/\y/\c in {0.35/0.30/cSem, 1.15/0.58/cEnc, 1.80/0.12/cSal}
      \node[star, star points=5, star point ratio=2.2, draw=\c, fill=\c!35,
            inner sep=0.7pt] at (\x,\y) {};
    \fill[black] (0.80,0.20) circle (0.042);
    \draw[dashed, black!55] (0.80,0.20) -- (0.35,0.30);
    \draw[dashed, black!55] (0.80,0.20) -- (1.15,0.58);
    \node[font=\tiny, text=black!60] at (0.52,0.02) {70\,\%};
    \node[font=\tiny, text=black!60] at (1.22,0.20) {30\,\%};}\\[6pt]
  \centering$s_{i,k}=\dfrac{e^{z_{i,k}}}{\sum_{j}e^{z_{i,j}}}$\\[6pt]
  \raggedright
  $s_{i,k}$ dice cuánto pertenece el punto $i$ al grupo $k$. Los valores de un
  punto suman 1, como repartir un porcentaje. $z_{i,k}$ es la puntuación que da
  la red y la exponencial solo sirve para convertirla en porcentaje.};

\node[tarj] (T3) at ($(pc)+(5.45,-1.15)$)
 {\textbf{\textcircled{\scriptsize 3} Los grupos se relacionan entre sí}\\[4pt]
  \centering
  \tikz[baseline=0, every node/.style={anchor=center}]{
    \fill[cSem!20, rotate=20] (0.72,0.26) ellipse (0.80 and 0.34);
    \fill[cEnc!20, rotate=-15] (1.42,0.52) ellipse (0.74 and 0.32);
    \foreach \x/\y in {0.25/0.20, 0.72/0.46, 1.05/0.14, 1.55/0.40, 1.98/0.12, 1.62/0.70}
      \fill[black!65] (\x,\y) circle (0.045);
    \node[font=\tiny, text=black!60] at (1.05,-0.18) {cada mancha es una hiperarista};}\\[6pt]
  \centering$S_b=\alpha\,S_g+(1-\alpha)\,S_s$\\[6pt]
  \raggedright
  Dos grupos se parecen por su forma ($S_g$) y por su significado ($S_s$).
  $\alpha$ reparte el peso: con $\alpha=1$ solo cuenta la forma y con
  $\alpha=0$ solo el significado. Cada grupo se enlaza con los $M=8$ más
  parecidos, y ese enlace de varios a la vez permite explicar una pieza nueva
  con varias conocidas.};

\draw[dashed, cFus!70, -{Latex[length=1.6mm]}]
  (M4.south) |- ($(T1.north)+(0,0.33)$) -- (T1.north);

%% ---------- de donde viene cada pieza ----------
\draw[-{Latex[length=1.8mm]}, semithick, black!55]
  ($(T1.north east)+(0,-2.85)$) -- ($(T2.north west)+(0,-2.85)$)
  node[midway, above=0.4mm, font=\tiny, text=black!60, align=center,
       fill=white, inner sep=1pt] {esa ficha\\entra aquí};
\draw[-{Latex[length=1.8mm]}, semithick, black!55]
  ($(T2.north east)+(0,-2.85)$) -- ($(T3.north west)+(0,-2.85)$)
  node[midway, above=0.4mm, font=\tiny, text=black!60, align=center,
       fill=white, inner sep=1pt] {los grupos\\ya formados};

\node[fit=(T1)(T2)(T3), inner sep=0pt, draw=none] (Tarjetas) {};
\node[anchor=north, font=\tiny, align=center, text=black!70, text width=13.2cm]
  at ($(Tarjetas.south)+(0,-0.3)$)
  {\textcolor{cSem}{$\blacksquare$}~Módulos 2 y 4: HyperNCD, el método que
   descubre clases nuevas.\quad
   \textcolor{cEnc}{$\blacksquare$}~Módulo 3: Sonata, red ya preentrenada que
   se usa sin volver a entrenarla.\quad
   \textcolor{cFus}{$\blacksquare$}~El reemplazo dentro del Módulo 4 es el
   aporte de esta tesis.};
\end{tikzpicture}
\end{document}
